\section{Consensus Taxonomy Summary}

\subsection{Winner by $k$}
\begin{table}[t]
\centering
\begin{tabular}{lrrr}
\hline
Slice & consensus\_all & sum\_minus\_emb & sum\_minus\_cc \\
\hline
top-$k$=6 & 52 & 15 & 0 \\
top-$k$=30 & 31 & 0 & 10 \\
Total & 83 & 15 & 10 \\
\hline
\end{tabular}
\caption{Bank-based local review winners by slice.}
\end{table}

\subsection{Taxonomy by winner}
\begin{table}[t]
\centering
\begin{tabular}{lrrrrr}
\hline
Taxonomy label & Count & Share & consensus\_all & sum\_minus\_emb & sum\_minus\_cc \\
\hline
broad\_context\_noise & 35 & 32.4\% & 34 & 1 & 0 \\
material\_family\_coherence & 29 & 26.9\% & 22 & 5 & 2 \\
method\_family\_coherence & 19 & 17.6\% & 15 & 4 & 0 \\
single\_cue\_specificity & 8 & 7.4\% & 4 & 1 & 3 \\
over\_regularized\_consensus & 6 & 5.6\% & 0 & 3 & 3 \\
application\_umbrella\_noise & 5 & 4.6\% & 2 & 1 & 2 \\
coherent\_refinement & 3 & 2.8\% & 3 & 0 & 0 \\
semantic\_drift & 3 & 2.8\% & 3 & 0 & 0 \\
\hline
\end{tabular}
\caption{Primary case taxonomy across all reviewed cases.}
\end{table}

\subsection{Representative cases}
\paragraph{consensus\_all}
\begin{itemize}
\item \textbf{The Catalytic Use of Onion-Like Carbon Materials for Styrene Synthesis by Oxidative Dehydrogenation of Ethylbenzene} (field\_15\_k30\_sum\_minus\_cc\_vs\_consensus\_bank\_n48, $k$=30, broad\_context\_noise, gap=4.0). Winner advantage: Group B provides an excellent, highly coherent local neighborhood. Nearly every paper discusses the target's specific topic: oxidative dehydrogenation using carbon-based catalysts, with several mentioning the exact ethylbenzene to styrene reaction.. Loser failure mode: Group A is extremely broad and noisy, only connected by the general theme of 'catalysis' and completely misses the specific catalyst family, reaction type, and chemical system of the target paper..
\item \textbf{Tunable thermomorphism and applications of ionic liquid analogues of Girard's reagents} (field\_15\_k30\_sum\_minus\_cc\_vs\_consensus\_bank\_n48, $k$=30, material\_family\_coherence, gap=3.5). Winner advantage: Group B provides an excellent, highly coherent neighborhood focused on the target's specific topic: using the temperature-dependent phase behavior (thermomorphism, LCST/UCST) of ionic liquids in water for extraction.. Loser failure mode: Group A is a very broad and noisy collection of papers on the general topic of ionic liquids, but it completely misses the target's core application and physical phenomenon..
\item \textbf{A New Category of Liquid Salt--Liquid Ionic Phosphates (LIPs)} (field\_15\_k30\_sum\_minus\_cc\_vs\_consensus\_bank\_n48, $k$=30, material\_family\_coherence, gap=3.5). Winner advantage: Group B provides an excellent local neighborhood. The top-ranked paper is a near-perfect match, discussing the same specific class of materials (LIPs from polyammonium salts). The subsequent papers create a highly coherent context by focusing on dicationic ionic liquids, which is directly relevant to the target's use of polycations. Group A, in contrast, is a very broad and noisy collection of general review papers and mismatched subtopics within the wider field of ionic liquids, failing to capture the specific innovations of the target paper.. Loser failure mode: Group B provides an excellent local neighborhood. The top-ranked paper is a near-perfect match, discussing the same specific class of materials (LIPs from polyammonium salts). The subsequent papers create a highly coherent context by focusing on dicationic ionic liquids, which is directly relevant to the target's use of polycations. Group A, in contrast, is a very broad and noisy collection of general review papers and mismatched subtopics within the wider field of ionic liquids, failing to capture the specific innovations of the target paper..
\item \textbf{On the solid, liquid and solution structural organization of imidazolium ionic liquids} (field\_15\_k30\_sum\_minus\_cc\_vs\_consensus\_bank\_n48, $k$=30, material\_family\_coherence, gap=3.5). Winner advantage: Group B presents an excellent, highly coherent neighborhood focused specifically on the structural organization, self-assembly, and properties of imidazolium ionic liquids, which is the direct topic of the target paper. Group A is significantly weaker and noisier; it contains very broad review-style papers, a news article, and a completely irrelevant paper on polymer drag reduction, all of which dilute the target's specific research context.. Loser failure mode: Group B presents an excellent, highly coherent neighborhood focused specifically on the structural organization, self-assembly, and properties of imidazolium ionic liquids, which is the direct topic of the target paper. Group A is significantly weaker and noisier; it contains very broad review-style papers, a news article, and a completely irrelevant paper on polymer drag reduction, all of which dilute the target's specific research context..
\end{itemize}
\paragraph{sum\_minus\_cc}
\begin{itemize}
\item \textbf{A single-site hydroxyapatite-bound zinc catalyst for highly efficient chemical fixation of carbon dioxide with epoxides} (field\_15\_k30\_sum\_minus\_cc\_vs\_consensus\_bank\_n48, $k$=30, single\_cue\_specificity, gap=2.5). Winner advantage: Group A provides an excellent and highly specific neighborhood. Almost every paper focuses on zinc-based catalysts for the coupling of CO2 and epoxides, which is the exact topic of the target paper. Group B is far too broad and noisy; while it generally addresses CO2 fixation, it uses a wide variety of different catalysts (Ni, Ru, Cr, salts) and completely misses the central 'zinc' theme of the target. Additionally, Group B includes irrelevant papers on polymerization.. Loser failure mode: Group A provides an excellent and highly specific neighborhood. Almost every paper focuses on zinc-based catalysts for the coupling of CO2 and epoxides, which is the exact topic of the target paper. Group B is far too broad and noisy; while it generally addresses CO2 fixation, it uses a wide variety of different catalysts (Ni, Ru, Cr, salts) and completely misses the central 'zinc' theme of the target. Additionally, Group B includes irrelevant papers on polymerization..
\item \textbf{Carbon Paste Electrode Modified with Functionalized Nanoporous Silica Gel as a New Sensor for Determination of Silver Ion} (field\_15\_k30\_sum\_minus\_cc\_vs\_consensus\_bank\_n48, $k$=30, single\_cue\_specificity, gap=2.0). Winner advantage: Group A provides a much more specific and coherent neighborhood for the target paper. It correctly identifies the two key research facets: the use of carbon paste electrodes (CPEs) for potentiometry and the specific application of detecting silver ions. Group B is plausible at a high level (potentiometric sensors) but is dominated by papers focused on determining organic drug molecules, which is a different sub-field from the target's focus on inorganic metal ions.. Loser failure mode: Group A provides a much more specific and coherent neighborhood for the target paper. It correctly identifies the two key research facets: the use of carbon paste electrodes (CPEs) for potentiometry and the specific application of detecting silver ions. Group B is plausible at a high level (potentiometric sensors) but is dominated by papers focused on determining organic drug molecules, which is a different sub-field from the target's focus on inorganic metal ions..
\item \textbf{Surface Analysis of Ionic Liquids with and without Lithium Salt Using X-ray Photoelectron Spectroscopy} (field\_15\_k30\_sum\_minus\_cc\_vs\_consensus\_bank\_n48, $k$=30, material\_family\_coherence, gap=2.0). Winner advantage: Group A provides an excellent, highly specific neighborhood. It includes multiple papers studying the exact same ionic liquid ([EMI][Tf2N]) with the same technique (XPS), and also captures the crucial context of adding lithium salts. Group B is on the correct general topic of surface analysis of ionic liquids with spectroscopy, but is much broader and noisier, with several papers focused on dissolved metal complexes rather than simple alkali salts, and it completely misses the lithium salt aspect of the target paper.. Loser failure mode: Group A provides an excellent, highly specific neighborhood. It includes multiple papers studying the exact same ionic liquid ([EMI][Tf2N]) with the same technique (XPS), and also captures the crucial context of adding lithium salts. Group B is on the correct general topic of surface analysis of ionic liquids with spectroscopy, but is much broader and noisier, with several papers focused on dissolved metal complexes rather than simple alkali salts, and it completely misses the lithium salt aspect of the target paper..
\item \textbf{Zn(<scp>ii</scp>)@TFP-DAQ COF: an efficient mesoporous catalyst for the synthesis of<i>N</i>-methylated amine and carbamate through chemical fixation of CO<sub>2</sub>} (field\_15\_k30\_sum\_minus\_cc\_vs\_consensus\_bank\_n48, $k$=30, single\_cue\_specificity, gap=2.0). Winner advantage: Group A provides a superb local neighborhood, with nearly every paper focused on the target's specific reactions: N-methylation, N-formylation, or carbamate synthesis using CO2. The catalyst theme is also very strong, featuring many Zn-based systems and porous frameworks like COFs and MOFs, including a near-perfect match at rank 5. Group B is much broader and noisier, centered on the general theme of CO2 fixation but with many papers targeting different products, making it a less precise description of the target's immediate research context.. Loser failure mode: Group A provides a superb local neighborhood, with nearly every paper focused on the target's specific reactions: N-methylation, N-formylation, or carbamate synthesis using CO2. The catalyst theme is also very strong, featuring many Zn-based systems and porous frameworks like COFs and MOFs, including a near-perfect match at rank 5. Group B is much broader and noisier, centered on the general theme of CO2 fixation but with many papers targeting different products, making it a less precise description of the target's immediate research context..
\end{itemize}
\paragraph{sum\_minus\_emb}
\begin{itemize}
\item \textbf{Impact of Ionic Liquid Physical Properties on Lipase Activity and Stability} (field\_15\_k06\_sum\_minus\_emb\_vs\_consensus\_bank\_n48, $k$=6, material\_family\_coherence, gap=2.5). Winner advantage: Group A provides a much better neighborhood by correctly identifying the target's core topic: lipase-catalyzed reactions in ionic liquids. Its top-ranked papers are highly relevant and coherent.. Loser failure mode: Group B is severely flawed, with its top-ranked papers focusing on an entirely different application of ionic liquids (metal extraction), demonstrating a clear failure to capture the target's immediate research context..
\item \textbf{Carbon dioxide induced separation of ionic liquids and water} (field\_15\_k06\_sum\_minus\_emb\_vs\_consensus\_bank\_n48, $k$=6, material\_family\_coherence, gap=2.5). Winner advantage: Group A provides an excellent, highly coherent neighborhood focused on the specific technique of using CO2 to induce phase separation in ionic liquid mixtures. While most papers discuss organic solvents, this is a very close parallel to the target's water-based system, and crucially, rank 8 is a direct match on the IL-water-CO2 system. Group B contains some of the same relevant papers but is severely degraded by high-ranking noise (two editorials) and a thematically distinct paper on CO2 capture, making it a much poorer representation of the target's immediate research context.. Loser failure mode: Group A provides an excellent, highly coherent neighborhood focused on the specific technique of using CO2 to induce phase separation in ionic liquid mixtures. While most papers discuss organic solvents, this is a very close parallel to the target's water-based system, and crucially, rank 8 is a direct match on the IL-water-CO2 system. Group B contains some of the same relevant papers but is severely degraded by high-ranking noise (two editorials) and a thematically distinct paper on CO2 capture, making it a much poorer representation of the target's immediate research context..
\item \textbf{Decoupling of ionic conductivity from structural dynamics in polymerized ionic liquids} (field\_15\_k06\_sum\_minus\_emb\_vs\_consensus\_bank\_n48, $k$=6, material\_family\_coherence, gap=2.0). Winner advantage: Group A provides an excellent local neighborhood by including multiple papers (ranks 4 and 7) that specifically address the core concept of "decoupling" of ionic conductivity from structural dynamics in polymerized ionic liquids. Group B, while topically relevant at a broader level, completely omits this key phenomenon. Therefore, Group A much more accurately captures the target paper's immediate and specific research context.. Loser failure mode: Group A provides an excellent local neighborhood by including multiple papers (ranks 4 and 7) that specifically address the core concept of "decoupling" of ionic conductivity from structural dynamics in polymerized ionic liquids. Group B, while topically relevant at a broader level, completely omits this key phenomenon. Therefore, Group A much more accurately captures the target paper's immediate and specific research context..
\item \textbf{Crystal polymorphism in 1-butyl-3-methylimidazolium halides: supporting ionic liquid formation by inhibition of crystallizationElectronic supplementary information (ESI) available: packing diagrams for I and II; table of closest contacts for I, I-Br and II. See http://www.rsc.org/suppdata/cc/b3/b304543a/} (field\_15\_k06\_sum\_minus\_emb\_vs\_consensus\_bank\_n48, $k$=6, material\_family\_coherence, gap=2.0). Winner advantage: Group A provides a superb local neighborhood. The top-ranked papers are all directly concerned with the crystal structure, polymorphism, and thermal properties of the exact ionic liquid family ([bmim] halides) as the target paper. Group B contains many of the same relevant papers, but its ranking is much weaker, placing a methods paper and a paper on a different ionic liquid too highly, which dilutes the focus and makes the neighborhood less coherent.. Loser failure mode: Group A provides a superb local neighborhood. The top-ranked papers are all directly concerned with the crystal structure, polymorphism, and thermal properties of the exact ionic liquid family ([bmim] halides) as the target paper. Group B contains many of the same relevant papers, but its ranking is much weaker, placing a methods paper and a paper on a different ionic liquid too highly, which dilutes the focus and makes the neighborhood less coherent..
\end{itemize}
